% One-time imported draft. Editable native source; no Markdown regeneration.
\section{One logical board, three physical
  representations}\label{008-life-generations-x-strides-one-logical-board-three-physical-representations}

{\def\LTcaptype{table} % do not increment counter
  \begin{longtable}[]{@{}
      >{\RaggedRight\arraybackslash}p{(\linewidth - 6\tabcolsep) * \real{0.2500}}
      >{\RaggedRight\arraybackslash}p{(\linewidth - 6\tabcolsep) * \real{0.2500}}
      >{\RaggedRight\arraybackslash}p{(\linewidth - 6\tabcolsep) * \real{0.2500}}
      >{\RaggedRight\arraybackslash}p{(\linewidth - 6\tabcolsep) * \real{0.2500}}@{}}
    \toprule\noalign{}
    \begin{minipage}[b]{\linewidth}\RaggedRight
      Allocation
    \end{minipage} & \begin{minipage}[b]{\linewidth}\RaggedRight
                       Cell width
                     \end{minipage} & \begin{minipage}[b]{\linewidth}\RaggedRight
                                        Physical row stride
                                      \end{minipage} & \begin{minipage}[b]{\linewidth}\RaggedRight
                                                         First logical cell
                                                       \end{minipage}                                                     \\
    \midrule\noalign{}
    \endhead
    \bottomrule\noalign{}
    \endlastfoot
    seed                                        & 1 byte                                      & 8 bytes                                     & seed+0  \\
    \texttt{board\_a} or \texttt{board\_b}      & 1 byte                                      & 10 bytes                                    & board
    base+11                                                                                                                                           \\
    frame                                       & 4 bytes                                     & 32 bytes                                    & frame+0 \\
  \end{longtable}
}

Interior logical offset is (y+1)\ensuremath{\times}10+(x+1) in a board, y\ensuremath{\times}8+x in seed, and
4\ensuremath{\times}(y\ensuremath{\times}8+x) in frame. Do not reuse one representation\textquotesingle s
stride for another.
